\documentclass[11pt]{article}
\usepackage{amsmath,amssymb,amsthm}
\newtheorem{lemma}{Lemma}\newtheorem{theorem}{Theorem}
\title{Records Never Stop: Infinitely Many Maximal Prime Gaps}
\author{Fang \& Fox Lemma Office (ShadowAlzazel, Kumi-San, Alo)}
\date{Status: \textbf{proven} (textbook; pipeline warmup)}
\begin{document}\maketitle

\paragraph{The picture.} Walk along the primes and keep a high-water mark of the
widest gap seen so far. Each time the water rises, that index is a \emph{record}.
The question is whether the water ever stops rising. It cannot: there are deserts
of composites as wide as we like, and the first gap that beats the old mark is,
by its very firstness, a new record.

\paragraph{Setup.} Write $p_n$ for the $n$-th prime ($p_0=2$) and $g_n=p_{n+1}-p_n$.
Index $n$ is a \emph{record} if $g_m<g_n$ for every $m<n$.

\begin{lemma}[Factorial desert \texttt{(gap\_unbounded)}]
For every $k\ge 1$ there is $n$ with $g_n\ge k+1$.
\end{lemma}
\begin{proof}
Let $N=(k+1)!$. For $2\le j\le k+1$, $j\mid N$ and $j\mid j$, so $j\mid N+j$ with
$N+j>j$; hence $N+2,\dots,N+k+1$ are all composite. Let $p_n$ be the largest prime
$\le N+1$ (it exists since $2\le N+1$). The next prime lies beyond the desert, so
$p_{n+1}\ge N+k+2$ and $g_n\ge k+1$.
\end{proof}

\begin{lemma}[First to beat the mark \texttt{(first\_exceeder\_is\_record)}]
Fix $M$ and let $H=\max_{m\le M} g_m$. If $n$ is the \emph{least} index with
$g_n>H$, then $n>M$ and $n$ is a record.
\end{lemma}
\begin{proof}
$n>M$ because every $g_m$ with $m\le M$ is $\le H$. Take $m<n$. If $m\le M$ then
$g_m\le H<g_n$. If $M<m<n$ then $g_m\le H$ by minimality of $n$, so again $g_m<g_n$.
\end{proof}

\begin{theorem}[\texttt{infinite\_maximalPrimeGapIndex}]
There are infinitely many record indices.
\end{theorem}
\begin{proof}
Given any $M$, Lemma~1 with $k=H$ gives some $g_n>H$; the least such $n$ is a record
beyond $M$ by Lemma~2. A set of naturals with an element past every bound is infinite.
\end{proof}

\paragraph{Why it matters for us.} Nothing here is deep; that is the point. The
\emph{shape} of the argument (an unbounded quantity forces infinitely many records,
via least-exceeder) is the same skeleton our \texttt{near\_misses.py} records use,
and it isolates exactly what the open Shanks record-index problem asks beyond this:
not \emph{whether} records recur, but \emph{how fast}, $g_n\sim(\log p_n)^2$.

\paragraph{Honesty.} Proven, elementary. No use of RH or any heuristic. The
Lean file should mirror Lemma~1, Lemma~2, Theorem in that order.
\end{document}
